% ═══════════════════════════════════════════════════════════════════
% INTRODUCTION GÉNÉRALE
% ═══════════════════════════════════════════════════════════════════

\chapter*{Introduction Générale}
\addcontentsline{toc}{chapter}{Introduction Générale}

Le rapprochement des écosystèmes de recherche entre l'Europe et l'Afrique s'appuie sur des collaborations scientifiques toujours plus denses : projets communs, échanges de chercheurs, réseaux d'experts et de laboratoires. La structuration de ces collaborations repose sur une ressource essentielle, les \textbf{contacts} : chercheurs, enseignants-chercheurs, industriels et partenaires institutionnels dont la mobilisation conditionne la réussite des projets de R\&I.

Or, au sein de l'ARSII, cette ressource était jusqu'alors gérée de façon manuelle --- fichiers tableurs éparpillés, sans référentiel commun ni traçabilité des mises à jour --- ou au moyen de solutions web payantes dont l'abonnement représente un budget non négligeable pour une association. Cette absence d'outil centralisé et maîtrisé constitue une véritable difficulté opérationnelle, l'association ne pouvant exploiter pleinement le potentiel de son réseau de contacts.

C'est dans ce cadre que s'inscrit ce stage d'été, dont l'objectif est de concevoir et réaliser une \textbf{plateforme web de gestion, d'organisation et d'enrichissement de bases de contacts} en R\&I Europe--Afrique. Ce rapport rend compte de l'ensemble du travail mené, de l'analyse des besoins à la validation finale de la solution.
